
%% $RCSfile: proj_proposal.tex,v $
%% $Revision: 1.4 $
%% $Date: 2017/10/06 02:55:50 $
%% $Author: kevin $

\documentclass[11pt, a4paper, twoside, openright]{article}

\usepackage{parskip}

\usepackage{float} % lets you have non-floating floats

\usepackage{url} % for typesetting urls

\usepackage[acronym]{glossaries} % for better handling of acronyms

%% Example acronyms used in this document. 
\newacronym{hci}{HCI}{Human-Computer Interaction}
\newacronym{acm}{ACM}{Association for Computing Machinery}
\newacronym{ieee}{IEEE}{Institute of Electrical and Electronics Engineers}

%% We don't want figures to float so we define
%%
\newfloat{fig}{thp}{lof}[section]
\floatname{fig}{Figure}

%% These are standard LaTeX definitions for the document
%%
\title{Scotland Yard Multiplayer Game as Map Overlay}
\author{Ming Bao}

%% This file can be used for creating a wide range of reports
%%  across various Schools
%%
%% Set up some things, mostly for the front page, for your specific document
%
% Current options are:
% [ecs|msor|sms]          Which school you are in.
%                         (msor option retained for reproducing old data)
% [bschonscomp|mcompsci]  Which degree you are doing
%                          You can also specify any other degree by name
%                          (see below)
% [font|image]            Use a font or an image for the VUW logo
%                          The font option will only work on ECS systems
%
\usepackage[image,ecs]{vuwproject} 

% You should specifiy your supervisor here with
%     \supervisor{Firstname Lastname}
% use \supervisors if there are more than one supervisor
\supervisors{Jens Dietrich, Stuart Marshall}

% Unless you've used the bschonscomp or mcompsci
%  options above use
%   \otherdegree{OTHER DEGREE OR DIPLOMA NAME}
% here to specify degree
\otherdegree{ENGR489}

% Comment this out if you want the date printed.
\date{}

\begin{document}
	
	% Make the page numbering roman, until after the contents, etc.
	\frontmatter
	
	%%%%%%%%%%%%%%%%%%%%%%%%%%%%%%%%%%%%%%%%%%%%%%%%%%%%%%%
	
	\begin{abstract}
		This project develops a web-based multiplayer version of Scotland Yard overlaid on a real-world map of the Wellington region. Physical gameplay limits accessibility, and existing digital versions lack local geographic integration. The system will implement core game logic and user interface components within a browser environment, integrating real-time multiplayer communication and interactive mapping technologies to support location-based gameplay. The expected outcome is a functional online platform that enables synchronous play, enforces game rules automatically, and visualises movement across the Wellington map. The project will be evaluated through functional testing of game mechanics, performance testing of multiplayer interactions, and user testing to assess usability and overall gameplay experience.
	\end{abstract}
	
	%%%%%%%%%%%%%%%%%%%%%%%%%%%%%%%%%%%%%%%%%%%%%%%%%%%%%%%
	\maketitle
	%\tableofcontents
	
	% we want a list of the figures we defined
	%\listof{fig}{Figures}
	
	%%%%%%%%%%%%%%%%%%%%%%%%%%%%%%%%%%%%%%%%%%%%%%%%%%%%%%%
	
	\mainmatter
	
	%\setlength{\parindent}{0pt}
	%\setlength{\parskip}{4pt}
	
	%%%%%%%%%%%%%%%%%%%%%%%%%%%%%%%%%%%%%%%%%%%%%%%%%%%%%%%
	
	\section{Introduction}
	
	Many board games are limited by physical components and the need for players to be in the same location. Scotland Yard, a hidden-movement strategy game based on movement across a map, is well suited to a digital implementation but most existing versions simply replicate the original board and do not use modern web or mapping technologies.
	
	This project aims to implement a web-based multiplayer version of Scotland Yard that overlays game logic and user interface elements onto a real-world map of the Wellington region using a map API. A browser-based frontend will provide an interactive interface for players, while a backend service will manage game state, enforce rules, and synchronise player actions in real time.
	
	The system will be evaluated through functional testing of the game mechanics, performance testing of multiplayer interactions, and user testing to assess usability and gameplay experience. The project will primarily require software resources including web development frameworks, mapping APIs, and networking tools for real-time communication.

	
	\section{The Problem}
	
	Board games such as Scotland Yard are traditionally limited to physical play, requiring co-located participants and a fixed board layout. Although digital versions exist, they typically replicate the original board and do not integrate real-world geographic data or modern web-based multiplayer capabilities. There is an opportunity to extend the game into a location-aware, browser-based platform that combines structured game logic with interactive mapping and contemporary user interface design.

	The engineering problem is to implement a web-based multiplayer system that enforces the rules of Scotland Yard while overlaying gameplay onto a real-world map of the Wellington region using an API from some map. This requires integrating turn-based game logic, real-time client–server communication, geographic data processing, and a responsive, intuitive frontend. The system must synchronise player actions reliably, ensure fairness through programmatic rule enforcement, and present clear visual feedback through effective interface design.
	
	The aim of this project is to implement and evaluate an online multiplayer version of Scotland Yard that runs entirely in a web browser, uses a map API to render Wellington as the game board, provides a well-designed frontend for interaction and visualisation, and supports real-time multiplayer gameplay. The project will demonstrate correctness of game mechanics, robustness of multiplayer synchronisation, and usability through structured testing and user evaluation.
	
	\section{Proposed Solution}
	
	In this section you will explain how you plan to solve the problem, that is, how you intend to carry the project out. At this early stage you need to be both clear about what you are going to do and flexible enough to adapt to changing circumstances. Making an early plan will not prevent you from running into trouble, but it will help you identify possible problems early. For example, if you intended to run an experiment in \gls{hci}, you might realise early on that there would be problems gathering sufficient data to get reliable results, and that you should re-design your experiment.
	
	Part of the planning process involves producing a timetable for when
	the work is actually going to be done.
	
	Each part of the project should produce some output. For example you
	might plan on spending two weeks on background reading: the output of
	this will be a bibliography, and a possibly a literature survey for
	your report. Indeed, if you take the advice given above about having
	specific, measurable goals, you should describe this part of your
	project as:
	
	\begin{description}
		\item[Good] Produce bibliography (est: 2 weeks)
	\end{description}
	rather than
	\begin{description}
		\item[\bf Bad] Background reading (est: 2 weeks)
	\end{description}
	
	Note that the methodology you outline here is dependent upon the type
	of project and engineering area. You must talk to your supervisor
	about this.
	
	\section{Evaluating your Solution}
	
	In this section you will explain how you will evaluate your solution
	once you have built it. The method of evaluation will be domain
	specific. Your supervisor should provide guidance as to what is an
	appropriate form of evaluation. For example, user testing for a \gls{hci} project or performance measurement for a Network Engineering project.
	
	\section{Resourcing and Ethics}
	
	In this section you will detail any resource requirements such as hardware, software or access. You should address each of these points, even if only to indicate that they are not of particular concern to your project.
	
	
	\subsection{Hardware}
	
	There aren't many hardware needs for this project. The main thing other than a laptop would be a server to host the project once its completed. I have a home server that I'm planning to host this on permanently once its finished so theres no extra costs.
	
	\subsection{Software, Datasets and Models}
	
	Discuss the specialised software (applications, compilers, libraries, development kits, simulators, solvers, etc) to be used. You need not include common desktop software such as browsers, office applications, IDEs, and editors. This is also the place to mention datasets and models that you plan to use. 
	
	\subsection{Space, Virtual Resources and Access}
	
	Discuss the space (for instance labs or specialised rooms, outdoor fields) and resources that can be accessed virtually (for instance computing grid, virtual machines, network testbeds) that your project will use, and how you plan to obtain access them. 
	
	\subsection{Budget}
	
	Provide an estimate of the budget for the project. If the budget is more than \$500 you should indicate any sources of funding outside the normal ENGR 489 budget. The purchase of special tools or software etc. may not need to be included in your \$500 limit, however, if something is necessary for your project you should definitely list it here.
	
	Do not pad the budget to get it up to \$500.
	
	My supervisor already ordered the physical boardgame which cost \$80
	I also want a claude code subscription for the year which cost \$400.
		
	
	\subsection{Ethics}
	
	No extra ethical considerations. Testing will done primarily through software and the user testing will be done among students.
	
	\subsection{Safety}
	
	The standard stuff to do with working with computers like eye strain, sitting for long periods of time and wrist pain.
	
	
	\subsection{Intellectual Property}
	
	IP will be mine but the university will have permission to use it for open day and advertisements.
	
\end{document}
